\documentclass[11pt]{article}
\usepackage[margin=1.1in]{geometry}
\usepackage{amsmath,amssymb}
\usepackage{booktabs}
\usepackage{graphicx}
\usepackage[hidelinks]{hyperref}
\usepackage{xcolor}

\title{A Shrinkage Law for Epistemic Transport at the Writing Step of
Language-Model Pipelines\\[4pt]
\large with an open measurement kit and two parameter cards}
\author{Sam Bobo\\ \small with analysis assistance from Claude (Anthropic)}
\date{Draft v0.1 --- August 2026}

\newcommand{\ww}{\ensuremath{w}}
\newcommand{\cc}{\ensuremath{c}}

\begin{document}
\maketitle

\begin{abstract}
When a language model writes the final answer of a multi-component pipeline,
the epistemic qualifications raised upstream---disclosed assumptions, staleness
warnings, jurisdictional distinctions, conditions of validity---do not
transport into the output. Across $1{,}340$ audited runs of two architectures
sharing one writer model, output qualification follows a single line:
$y = \ww s + \cc$, the form of a shrinkage estimator that weighs upstream
supply $s$ against a prior. Under a conditional-preservation instruction the
evidence weight is $\ww=0.35$ $[0.31,0.39]$; under the canonical
mixture-of-agents prompt it is $\ww=0.16$ $[0.08,0.23]$, statistically
inseparable from ignoring upstream entirely; and compensating invention at
zero supply---qualification fabricated when no component raised any---appears
in both architectures ($40\%$ and $28\%$ of zero-supply runs). Three further
results sharpen the law: the parameters are set by the writing instruction,
not the architecture; corroboration of a claim by multiple components does not
raise its transport on our validated measurement channel; and no intervention
we tested (four instructions, three training objectives, one correction loop)
moved the parameters durably, while post-hoc selection with a verifier the
generator never observes reliably does. We release a zero-dependency
measurement kit that estimates the parameters from any pipeline's logs,
carrying the identifiability guards that disqualified $21$ of our own $27$
experimental arms and the instrument-validation protocol without which our own
second instrument measured nothing. Finally we derive the \emph{optimal}
parameters from upstream reliability, showing that perfect transport
($\ww{=}1$, $\cc{=}0$) is optimal only for perfectly reliable components: the
parameters are diagnostics for choosing interventions, never objectives to
train or search against.
\end{abstract}

\section{Introduction}
\label{sec:intro}

Multi-component language-model systems end, almost universally, in a writing
step: one model reads what specialists, proposers, or retrieved documents
produced and writes the answer a user sees. A large literature measures
whether this step preserves \emph{accuracy}. This paper measures whether it
preserves \emph{epistemic content}---the assumptions a component disclosed,
the staleness it warned about, the conditions under which its claim holds.
These are the parts of an answer that tell a reader how much to trust the
rest, and they are precisely the parts a downstream reader cannot
reconstruct if the writer drops them.

Our central finding is that the writing step does not transport this content.
It behaves as an estimator: it emits the level of qualification an answer of
its type ought to carry, weighing the upstream evidence against a prior
learned long before the pipeline existed. That behavior is captured by one
measurable law, and the law by two parameters,
\begin{equation}
\label{eq:law}
y \;=\; \ww\, s \;+\; \cc ,
\end{equation}
where $s$ is the number of qualification families the upstream components
raised, $y$ the number in the final output, $\ww$ the \emph{evidence weight},
and $\cc$ the \emph{prior fill}. A transducer has $\ww=1$, $\cc=0$. A writer
that ignores its inputs has $\ww=0$. Every system we measured sits low in
between, and the two corollaries of sitting there are the phenomena a
deployer actually experiences: at high supply the writer discards most of
what its components raised, and at zero supply it \emph{invents}---emitting
qualifications with no upstream source, at rates we measure at $28\%$
and $40\%$ of zero-supply runs in the two architectures.

The contributions are:

\begin{enumerate}
\item \textbf{The law, measured} (\S\ref{sec:cards}): parameter cards for two
architectures sharing one writer---a routed specialist council and a
canonical mixture-of-agents (MoA) stack---with bootstrap intervals, a
five-level supply sweep in which zero supply is measured rather than
extrapolated, and per-family decomposition on a validated channel.
\item \textbf{What sets the parameters} (\S\ref{sec:sets}): the writing
instruction moves \ww{} more than the entire architectural difference
($0.16 \to 0.35$); the compensation phenomenon travels across architectures;
and corroboration---the one reliability signal a multi-agent system uniquely
produces---does not modulate transport on the validated channel.
\item \textbf{Interventions under the law} (\S\ref{sec:interventions}):
instructions, preference training, and detector-fed correction loops all
fail, the last by producing improvement only its own instrument can see;
verifier-blind selection works, and its yield must be computed from the
\emph{measured} best-of-$n$ curve, because redraws are positively correlated
and the independence formula over-promises in both architectures.
\item \textbf{The kit} (\S\ref{sec:kit}): a zero-dependency package that
estimates $(\ww,\cc)$, the feature-span fraction, instrument blind spots, and
the clean-sample rate from logs, with identifiability guards and a
per-channel instrument-validation protocol, both of which materially changed
our own conclusions.
\item \textbf{The optimal parameters, derived} (\S\ref{sec:optimal}):
$\ww^{*}$ and $\cc^{*}$ follow from measured upstream reliability and
declared error costs; $\ww^{*}{=}1$ requires perfectly reliable components,
and $\cc^{*}{>}0$ is optimal whenever components miss warranted caution. The
parameters are therefore diagnostics, never training or search objectives.
\end{enumerate}

Scope, stated before anything else: both architectures were built in one lab,
share one writer model and one authored case battery, and all rates are
relative to a lexicon whose per-family validity we measured and report. The
paper's claims are calibrated to that evidence, its falsification conditions
are explicit (\S\ref{sec:falsify}), and the kit exists so that any reader can
produce the third card.

\section{Definitions and estimators}
\label{sec:defs}

\textbf{Property class.} We measure four families of epistemic
qualification: disclosure of training-data staleness (\emph{cutoff}),
labeling of numbers as assumptions (\emph{modeled}), separate treatment of
jurisdictions or regimes (\emph{jurisd}), and statements of conditions under
which a claim changes (\emph{hedging}). The framework is not specific to
this class; the kit ships a second class (source attribution) as a plug-in.

\textbf{Supply and emission.} For one run, $s$ is the number of families
present in the text the writer actually received (verbatim capture is a
precondition; see \S\ref{sec:kit}), and $y$ the number present in its
output. Runs shorter than $500$ characters are excluded, never scored.

\textbf{The regression.} $(\ww,\cc)$ are the OLS coefficients of
Eq.~\ref{eq:law} over runs, with seeded bootstrap intervals. Because $y$
counts at most four families, the fit is reported alongside its
identifiability guards: the estimate is refused when supply spans fewer than
three levels, flagged as \emph{extrapolated} when no runs at $s{=}0$ exist
(an intercept fitted outside its data routinely returns the impossible
$\cc<0$), and flagged as \emph{weakly identified} when the \ww{} interval
exceeds $0.35$. Applied to our own ledger, these guards disqualified $21$ of
$27$ arms with $n\geq30$---most of our own experiments were never powered
for the comparisons their point estimates invited.

\textbf{Per-family transmission.} Family presence is binary per run, so
per family the regression is unidentifiable and the registered quantities
are conditional: $T_f = P(f \in \text{output} \mid f\ \text{raised})$,
$I_f = P(f \in \text{output} \mid f\ \text{not raised})$, and the
discrimination $D_f = T_f - I_f$.

\textbf{Auxiliary parameters.} The feature-span fraction $f$ (share of
output characters in qualification-bearing sentences) quantifies why
sequence-level preference training dilutes; the clean-sample rate $p_0$
(probability a run contains no invented family) sizes selection
(\S\ref{sec:interventions}).

\section{Instruments, validated per channel}
\label{sec:instruments}

Every number in this paper is instrument-relative, so the instruments were
graded before the numbers were trusted. Against blinded dual-judge labels
over $200$ sentences (judge--judge agreement $0.86$; accuracy $0.90$/$0.93$
on hand-written anchors that include paraphrase positives and keyword-bait
negatives), the lexicon detects the \emph{modeled} family well (sensitivity
$0.92$, precision $0.92$) and the remaining families conservatively
(specificity $\geq 0.95$, sensitivity $0.25$--$0.30$). All composite rates
are therefore undercounts outside the modeled family, and \emph{modeled} is
the validated channel on which per-family claims rest.

The cautionary result: our second instrument, a natural-language-inference
detector with a genuine $0.93$ AUC on the discrimination task it was
calibrated for, was re-calibrated against the same judge labels for the task
it had actually been used on---family presence---and carried \emph{no
signal}: per-family AUC $0.12$--$0.55$, two families inversely related to
truth. No threshold rescues an AUC below chance. Two instruments had
"agreed" for weeks while measuring different constructs. The protocol that
found this ships with the kit, and the framework's rule follows from it:
a verdict requires two instruments \emph{whose sensitivity to the construct
has itself been demonstrated}; agreement between unvalidated instruments is
rhetoric.

\section{Two parameter cards}
\label{sec:cards}

\begin{table}[t]
\centering
\small
\begin{tabular}{lcc}
\toprule
 & \textbf{Council} & \textbf{Mixture-of-Agents} \\
 & (routed specialists, & (general proposers, \\
 & conditional instruction) & canonical prompt) \\
\midrule
runs & $1{,}260$ & $80$ \\
evidence weight \ww & $0.352\ [0.310, 0.391]$ & $0.158\ [0.076, 0.234]$ \\
prior fill \cc & $0.540\ [0.437, 0.651]$ & $0.275\ [0.104, 0.456]$ \\
prior-trust ratio $(1{-}\ww)/\ww$ & $1.84$ & $5.35$ \\
supply strata means ($s{=}0..4$) & $.40\,/\,.81\,/\,1.32\,/\,1.62\,/\,1.83$
 & $.28\,/\,.44\,/\,.50\,/\,.88\,/\,.83$ \\
strata $R^2$ / run-level $R^2$ & $0.976$ / $0.140$ & $0.890$ / $0.136$ \\
zero-supply runs inventing & $40\%\ [29\%, 52\%]$ & $27.8\%\ [12\%, 51\%]$ \\
feature-span fraction $f$ & $0.068$ & $0.002$ \\
$D_{\text{modeled}}$ (validated channel) & $0.213\ [0.155, 0.271]$
 & $-0.010\ [-0.139, 0.117]$ \\
\bottomrule
\end{tabular}
\caption{Parameter cards for two architectures sharing one writer model
(gpt-oss-20B). Both $R^2$ values are reported deliberately: the law
describes the conditional mean tightly and individual runs loosely, and
quoting only the strata fit would oversell it. The council card pools all
usable ledger runs; the MoA card is an $80$-run supply sweep in which supply
was manipulated by sentence-level ablation and injection of the upstream
texts, so $s{=}0$ is measured, not extrapolated.}
\label{tab:cards}
\end{table}

Table~\ref{tab:cards} is the paper's central evidence. Three readings:

\textbf{The form travels.} Both architectures show monotone strata, positive
prior fill, and an evidence weight whose interval excludes zero: writers
respond to supply, incompletely, and fill from a prior when supply is
absent. Compensating invention at zero supply---the deployer-facing
symptom---appears in both ($40\%$ and $28\%$ of zero-supply runs), and on
trigger-light cases authored to warrant \emph{no} qualification, so for
those cases invention is miscalibration by construction, independent of any
judgment about upstream quality.

\textbf{The parameters do not travel.} The MoA card sits near the
pure-register corner: its \ww{} interval does not exclude $0.15$, its
prior-trust ratio is three times the council's, and on the validated channel
its discrimination is indistinguishable from zero---the aggregator emits
modeled qualification at the same rate whether or not any proposer raised
it. Its composite slope is carried by families whose distinctive phrases
the writer sometimes echoes, not by behavior on the channel we can measure
cleanly. We registered the prediction that the MoA card would show the full
shrinkage form with both corners excluded; it did not, and per the
registration this boundary is reported at equal prominence with the
successes.

\textbf{The near-silent regime exists and scores deceptively well.} The MoA
outputs carry almost no qualification content at all ($f=0.002$; the
council's naive flat-merge arm measured $f=0.000$). A silent writer invents
little because it says little---composite invention metrics reward it, and
any evaluation that does not report $f$ beside the invention rate will
mistake silence for calibration.

\section{What sets the parameters}
\label{sec:sets}

\textbf{The instruction, more than the architecture.} Holding the writer
model fixed and restricting to a shared nine-case battery: the canonical MoA
prompt over general proposers yields $\ww=0.16\ [0.08,0.23]$; a naive merge
prompt over domain specialists yields $\ww=0.255\ [0.13,0.39]$; the
conditional-preservation instruction over the same specialists yields
$\ww=0.346\ [0.22,0.47]$. The instruction change moves \ww{} by more than
the upstream-composition change, and no instruction we found moves it past
roughly a third of the way to transport. Clause ablation in the earlier
program localizes the effect further: a single conditional clause carries
essentially the entire instruction effect, and each clause lifts primarily
the family it names. Instructions act as a gain control on transport---
graded, targeted, and bounded well short of fidelity.

\textbf{Corroboration is not used.} A multi-agent system produces one
reliability signal a single model cannot: independent agreement among
components. If the writer performed any reliability weighting, a family
raised by two or more components should transport at a higher rate than a
family raised by one. On the validated channel it does not:
$T(k{=}1)=0.524$ ($n{=}824$) versus $T(k{\geq}2)=0.487$ ($n{=}76$),
difference CI $[-0.153, +0.080]$. The three low-validity families show
large monotone agreement effects (up to $+0.43$), which we report without a
mechanism claim: genuine family-specific weighting, phrase-exposure
artifact (more components using a distinctive phrase mechanically raises
the chance one detectable phrasing survives), and case-mix confounding are
all live, and separating them requires judge labels rather than a lexicon.
On the channel we can trust, the summary of this section is one sentence:
\emph{instructions set the gain; nothing we found sets the calibration}---
not architecture, not training (\S\ref{sec:interventions}), and not the
architecture's own agreement structure.

\section{Interventions under the law}
\label{sec:interventions}

The law explains why the obvious interventions fail, and the failures in
turn support the law. We summarize; the companion behavioral paper reports
each in full.

\textbf{Instructions} move surface density (the gain of
\S\ref{sec:sets}) but leave invention intact: an explicit
suppression clause left zero-supply invention statistically unchanged.
Under the law this is expected---shrinkage toward a prior is loss-optimal
on the training distribution, and no instruction changes what is
loss-optimal.

\textbf{Sequence-level preference training} (DPO/ORPO variants, three
objective classes) produced the dissociation the feature-span fraction
predicts: preference accuracy rose $0.48 \to 0.94$ while emitted behavior
did not move. With $f \approx 0.07$, roughly $93\%$ of a sequence-level
margin is satisfiable off-feature; we measured the pair construction used
in the failed runs at $0.87$ off-feature mass, against $0.00$ for
programmatically constructed minimal pairs.

\textbf{Detector-fed correction} is worse than nothing. A runtime audit
that quoted its detector's matched phrases back to the writer removed
$88\%$ of flagged inventions \emph{by that detector's count}; manual
inspection showed the qualifications surviving in rewording, and one
flagged "invention" was the detector's own false positive. The instrument
that generates feedback can never grade compliance---and a second
instrument only helps if its sensitivity is demonstrated
(\S\ref{sec:instruments}).

\textbf{Verifier-blind selection} is the intervention the law permits.
Sample $n$ outputs, keep the one whose provenance audit is cleanest, never
reveal the verifier to the generator; annotate residual failures rather
than revising them. Structurally, selection conditions the output
distribution while feedback conditions the policy input; only the latter
can be gamed, because only the latter is observable. Sizing must use the
\emph{measured} curve: in both architectures the empirical best-of-$n$
yield sits below the independence formula $1-(1-p_0)^n$ (council:
$0.83$ measured vs.\ $0.94$ predicted at $n{=}3$; MoA: $0.78$ vs.\ $0.89$
at $n{=}2$), because a writer that invents on a prompt tends to invent
again on redraws. The kit computes both curves and reports how many
distinct tasks support each point.

\section{The measurement kit}
\label{sec:kit}

The kit (\texttt{gst-kit}, MIT, zero dependencies in the core) turns a
pipeline's logs into a parameter card. Integration is one adapter from the
system's log shape to a four-field record---upstream texts as the writer
received them, output, task identifier, condition---with prebuilt adapters
for MoA stacks, LangGraph- and AutoGen-style frameworks, retrieval
pipelines, and generic JSONL. Three design commitments came from
measurement failures documented in our own program:

\begin{enumerate}
\item \textbf{Guards over estimates.} Every estimator refuses or flags an
answer its input could not identify (supply that never varies, an
extrapolated intercept, a weakly identified slope, best-of-$n$ points
supported by an unrepresentative subpopulation). Each guard exists because
its absence produced a wrong number that we published internally or nearly
published.
\item \textbf{Execution-path assertion.} The writer prompt actually used is
captured and asserted per run; violating runs are quarantined and counted.
Thirty-nine of our runs silently measured a prompt that never executed, and
two verdicts had to be withdrawn.
\item \textbf{Structural safety over procedural care.} The selection API
takes a zero-argument sampler; passing a callable that could receive
feedback is a type error, not a guideline.
\end{enumerate}

\section{The optimal parameters are derived, not chosen}
\label{sec:optimal}

It is tempting to read Table~\ref{tab:cards} as a scoreboard on which
$\ww{=}1$, $\cc{=}0$ wins. It is not, and the reason disciplines the whole
framework. The writer faces a signal-detection problem: upstream supply is
a noisy signal of whether qualification is \emph{warranted}. The optimal
response to a noisy signal is exactly a shrinkage estimator, with
\begin{equation}
\ww^{*} \;=\; \frac{\tau^{2}}{\tau^{2}+\sigma^{2}},
\end{equation}
the upstream components' signal-to-noise ratio. Perfect transport is
optimal only for perfectly reliable components. Our specialists are not:
they raise qualification on content authored to warrant none, so
$\ww^{*}<1$; and to the extent they \emph{miss} warranted caution,
$\cc^{*}>0$---some invention is a writer correctly filling gaps. The
measured pathology is therefore not shrinkage itself but
\emph{miscalibrated} shrinkage: weights set by pretraining priors rather
than by the measured reliability of the actual upstream, and invention
uncorrelated with warrant (the zero-supply inventions occur on cases
authored to warrant nothing).

Two consequences. First, the honest scope of our normative claims:
"the writer under-weights its components" presumes the components deserve
the weight, and upstream warrant-precision is a quantity we have not
measured; our unconditional claims are anchored on the authored zero-warrant
cases. Second, and centrally: because $(\ww^{*},\cc^{*})$ are functions of
measured upstream reliability and \emph{declared} error costs, legitimate
optimization targets the derivation---never the diagnostic. Training or
searching against $(\ww,\cc)$ directly selects for transcription and for
instrument artifacts (a perfect score describes a parrot, not a judge), and
at ecosystem scale a reported diagnostic becomes a target without anyone
optimizing it deliberately. The framework's rule: measure $(\ww,\cc)$ to
choose an intervention; estimate $(\ww^{*},\cc^{*})$ from held-out
reliability measurements and stated costs if a target is needed; and keep
evaluation batteries sealed from tuning.

\section{Related work}
\label{sec:related}

Mixture-of-agents aggregation was introduced by Wang et
al.~\cite{wang2024moa}; Li et al.~\cite{li2025selfmoa} showed aggregation
quality is sensitive to proposer quality, the accuracy-domain analogue of
our reliability argument. A 2026 literature documents uncertainty loss in
agent pipelines: Shi et al.~\cite{shi2026laundering} name the phenomenon of
fragile upstream states repackaged as confident artifacts, and propagation
frameworks~\cite{up2026,runtime2026,confsignals2026} monitor confidence
across components. Relative to this line, our contributions are (i)
\emph{bidirectionality}---we measure invention under scarcity, not only
loss; (ii) a \emph{quantitative law} with estimable parameters rather than
a qualitative failure mode; (iii) an \emph{operational calibration test}
for aggregators (fitted \ww{} against derived $\ww^{*}$); and (iv) the
instrument-validation protocol without which, our own experience shows,
such measurements are unreliable. Single-model neighbors measure certainty
distortion between input and output~\cite{cd2026} and uncertainty
preservation in clinical summarization~\cite{clinical2026}; neither
addresses the multi-component writing step or supplies portable
measurement. The learning--acting dissociation under sequence-level
preference objectives connects to known credit-assignment limits of
DPO-family losses~\cite{rafailov2023dpo,hong2024orpo}; Goodhart dynamics
under optimization pressure are classical~\cite{goodhart1984,campbell1979}
and we treat our detector-fed correction result as a measured instance.

\section{Limitations}
\label{sec:limits}

Both cards come from one lab, one writer model, one authored nine-case
battery, and one property class; the MoA card is corroboration, not
independent replication. All rates are lexicon-relative; outside the
modeled family the lexicon is a high-precision undercounter, and per-family
claims are made only on the validated channel. Supply in the MoA sweep was
manipulated by ablation and injection, so its high-supply strata measure
mixed upstream. Upstream warrant-precision is unmeasured, which bounds the
normative reading of the discard results (\S\ref{sec:optimal}); the
invention results are anchored on authored zero-warrant cases and do not
share that bound. $\ww$ and $D_f$ are different estimands; we compare them
only where registered. The agreement-effect mechanism on low-validity
families is unresolved. All experiments were pre-registered with
consequences stated in advance; registrations, amendments, and two
corrections issued during this program's own audit are in the public
runbook.

\section{Falsification and replication}
\label{sec:falsify}

The framework is refuted, not embarrassed, by: a writer trained on ordinary
text showing $\ww\approx1$, $\cc\approx0$ without intervention; a supply
response no monotone estimator reading survives; sequence-level preference
training moving a small-$f$ behavior at modest scale; or selection showing
improvement under its own verifier that an independent validated instrument
denies. The cleanest single test is the selection-versus-feedback
asymmetry: the same verifier, used post-hoc versus in-loop, is predicted to
produce real versus instrument-relative improvement. Running any of these
requires an afternoon: one adapter, one \texttt{gst measure} invocation,
and a supply sweep by ablation if natural variation is insufficient. The
third parameter card---anyone's---is the most useful thing this paper can
cause to exist.

\section{Conclusion}
\label{sec:conclusion}

One line describes what a pipeline's writing step does to epistemic
content; two numbers describe how badly; both are measurable from logs in
an afternoon. In every configuration we built, the writer weighed its own
prior above its components' evidence, filled silence with invented caution
on cases authored to warrant none, ignored the one reliability signal its
architecture uniquely produced, and resisted every correction except being
sampled and selected by a verifier it never saw. The instruction is the
only dial that moves the evidence weight, and it stops a third of the way
to fidelity. Whether these parameters are constants of current models or
artifacts of our corner of the design space is exactly the question the
kit exists to settle.

\begin{thebibliography}{9}
\bibitem{wang2024moa} J.~Wang et al. Mixture-of-Agents enhances large
language model capabilities. \emph{arXiv:2406.04692}, 2024.
\bibitem{li2025selfmoa} W.~Li et al. Rethinking Mixture-of-Agents: is
mixing different large language models beneficial?
\emph{arXiv:2502.00674}, 2025.
\bibitem{shi2026laundering} K.~Shi et al. Confidence laundering in agent
systems: why uncertainty needs a latent carrier. \emph{arXiv:2606.20662},
2026.
\bibitem{up2026} Uncertainty propagation in LLM-based systems.
\emph{arXiv:2604.23505}, 2026.
\bibitem{runtime2026} Runtime uncertainty monitoring for LLM-based
multi-agent systems using Bayesian networks. \emph{arXiv:2607.25877}, 2026.
\bibitem{confsignals2026} Multiagent protocols with aggregated confidence
signals. \emph{arXiv:2606.13591}, 2026.
\bibitem{cd2026} Certainty distortion in large language models.
\emph{arXiv:2606.07951}, 2026.
\bibitem{clinical2026} Do language models preserve clinical uncertainty? A
benchmark for uncertainty-aware summarization. \emph{arXiv:2606.18471},
2026.
\bibitem{rafailov2023dpo} R.~Rafailov et al. Direct preference
optimization: your language model is secretly a reward model.
\emph{NeurIPS}, 2023.
\bibitem{hong2024orpo} J.~Hong, N.~Lee, J.~Thorne. ORPO: monolithic
preference optimization without reference model. \emph{arXiv:2403.07691},
2024.
\bibitem{goodhart1984} C.~Goodhart. Problems of monetary management: the
UK experience. In \emph{Monetary Theory and Practice}, 1984.
\bibitem{campbell1979} D.~Campbell. Assessing the impact of planned social
change. \emph{Evaluation and Program Planning}, 2(1), 1979.
\end{thebibliography}

\end{document}
